\par\noindent\SourceLeaf{57}%
\hspace*{\parindent}Die andere hingegen

\[
\begin{aligned}
&n\cos\varphi\bigl(A^{(1)}\sin\varphi+2A^{(2)}\sin2\varphi+.\,.\,.\\&\quad+hA^{(h)}\sin h\varphi+(h+1)A^{(h+1)}\sin(h+1)\varphi+.\,.\,.\bigr),
\end{aligned}
\]

\noindent enthält in jedem einzelnen Glied das Produkt aus $\cos\varphi$ und dem Sinus eines Vielfachen von $\varphi$. Um sie mit der vorigen zu einer einzigen Summe zu vereinigen, muss man sie mithilfe der bekannten Formel

\[
\cos\varphi\sin\lambda\varphi=\frac{1}{2}\sin(\lambda+1)\varphi+\frac{1}{2}\sin(\lambda-1)\varphi,
\]

\noindent umformen. So geht das Glied, das in ihr $\sin\varphi$ enthält, aus dem Glied

\[
n\cos\varphi.2A^{(2)}\sin2\varphi
\]

\noindent hervor und wird zu $n.A^{(2)}\sin\varphi$. Sucht man in ihr allgemein die Glieder, die $\sin(h+1)\varphi$ enthalten, so muss man das Glied

\[
hA^{(h)}\sin h\varphi.n\cos\varphi
\]

\noindent (und aus seiner Entwicklung den ersten Teil

\[
\frac{nhA^{(h)}}{2}\sin(h+1)\varphi)
\]

\noindent zusammen mit dem Glied

\[
(h+2)A^{(h+2)}\sin(h+2)\varphi.n\cos\varphi
\]

\noindent (und aus dessen Entwicklung den zweiten Teil

\[
\frac{n(h+2)A^{(h+2)}}{2}\sin(h+1)\varphi);
\]

\noindent nehmen. Damit erhält man als Koeffizienten von $\sin(h+1)\varphi$ in der gesuchten Reihe

\[
\frac{nhA^{(h)}}{2}+\frac{n(h+2)A^{(h+2)}}{2}.
\]

$\gamma$) Da nun die Reihe auf der einen Seite ($\alpha$) gleich den Reihen auf der anderen Seite ($\beta$) ist, müssen die entsprechenden Koeffizienten gleich sein, also hier diejenigen, die mit dem Sinus desselben Winkels multipliziert werden. So ist auf der einen Seite der Koeffizient von $\sin\varphi$

\[
=\nu nA-\nu\frac{nA^{(2)}}{2}
\]

\noindent auf der anderen Seite aber $A^{(1)}+nA^{(2)}$. Setzt man beide gleich, so erhält man

\[
\nu nA-\nu\frac{nA^{(2)}}{2}=A^{(1)}+nA^{(2)},\qquad\mbox{also}\qquad A^{(2)}=\frac{2(\nu nA-A^{(1)})}{(\nu+2)n}.
\]

